% ── SECTION 8: CONCLUSION (~600 words) ───────────────────────────────────────

\section{Conclusion: What Mutual Validation Means}
\label{sec:conclusion}

% Key points:
%   - What it means that three traditions arrived independently at the same structure.
%   - Not proof of any single tradition's authority.
%   - Evidence that all three are detecting a real feature of reality.
%   - Overlapping consensus: each tradition validates the others without dissolving.
%   - Implications for philosophy of science, science-religion dialogue, consciousness.

Three intellectual traditions --- quantum mechanics, Whitehead's process
philosophy, and the \bahai writings --- developed in substantially
independent contexts and using substantially different methods, converge
on a structurally identifiable description of reality: existence as
constituted by relationship, properties as emergent through relational
interaction, and a single relational principle --- named \mahabbat in
the \bahai lexicon, \emph{Eros} and creative advance in Whitehead, and
the nonlocal coherence of entanglement in quantum theory --- that
performs, in each vocabulary, the same structural work of sustaining
integrated being.

The process ontology bridge transforms this observation into an argument.
Quantum mechanics does not merely resemble process ontology: every
interpretation of quantum mechanics that survives the experimental
record (Bell~\citeyear{Bell1964} and its successors) either is
explicitly a process ontology or collapses, on examination, into a
relabeled version of one.
The \bahai writings do not merely anticipate process ontology; they
instantiate it, decades before \citet{Whitehead1929} named the framework
philosophically and more than half a century before Bell's theorem
formalized the experimental constraint that ruled its classical local
alternative out.
Process philosophy does not merely mediate between physics and spirituality;
it names the structure that both independently discovered.

The overlapping consensus architecture preserves the integrity of each tradition.
No tradition is asked to dissolve into the others or to accept the others'
methodology as authoritative.
Each tradition validates its own thesis on its own grounds.
The convergence is the finding --- and it is more powerful than any single tradition
could produce alone.

What the convergence means for the future of science-religion dialogue, for the
philosophy of science, and for our understanding of what reality is made of, is
a question this paper opens rather than closes.
That is, perhaps, the most it can do.
And it may be enough.
